\section{The global Hall--skein subquotient theorem}
\label{sec:subquotient}

We can now prove the global algebra statement.  All compatible products were
already correct before quotienting; the ideal $\Jsk$ supplies precisely the
missing crossing relations.

\subsection{The canonical map from the skein presentation}

By \cref{cor:arc-generation}, the quotient $\Hred/\Jsk$ is generated by the
classes $\overline\eps_\gamma$ of diagonals and by the frozen boundary units.
The defining Hall--skein defects give
\begin{equation}
 \overline\eps_\gamma\,\overline\eps_\delta
 =\omega^{a(\gamma,\delta)}
   \overline\eps_{\gamma\#_0\delta}
 +\omega^{b(\gamma,\delta)}
   \overline\eps_{\gamma\#_\infty\delta}
 \label{eq:quotient-skein-relation}
\end{equation}
for every ordered crossing pair.  Compatible arc classes obey the required
quasi-commutation relation by
\eqref{eq:compatible-Hall-quantum}; boundary and contractible relations were
imposed in \cref{def:Hred}.  The polygon skein presentation
\eqref{eq:skein-presentation} therefore gives a surjective map
\begin{equation}
 \Xi:
 \Skein_\omega^{\rm fr}(P)
 \longrightarrow
 \Hred/\Jsk,
 \qquad
 s_\gamma\longmapsto\overline\eps_\gamma.
 \label{eq:skein-to-Hall-quotient}
\end{equation}

\begin{theorem}[Hall--skein subquotient realization]
\label{thm:Hall-skein-subquotient}
Let $A$ be a cluster-tilted algebra of Dynkin type $A_n$ over
$\mathbb F_q$, let $P=P_{n+3}$ be the polygon model, and fix a
Hall-compatible quantum coefficient system.  Then the map
\eqref{eq:skein-to-Hall-quotient} is an isomorphism.  Consequently there is a
surjective algebra homomorphism
\begin{equation}
 \Phi_{\rm sk}:
 \Hred\longrightarrow\Aq
 \label{eq:global-Hall-cluster-map}
\end{equation}
with
\begin{equation}
 \Ker\Phi_{\rm sk}=\Jsk,
 \qquad
 \Hred/\Jsk\xrightarrow{\sim}\Aq.
 \label{eq:global-subquotient}
\end{equation}
For every diagonal $\gamma$,
\[
 \Phi_{\rm sk}(\eps_\gamma)=X_\gamma,
\]
and for every non-crossing multidiagonal $D$,
\[
 \Phi_{\rm sk}(\eps_D)
 =M_R(\mult_D)=\Theta_{g(D)}.
\]
\end{theorem}

\begin{proof}
Surjectivity of $\Xi$ was established above.  By
\cref{thm:skein-basis}, the source has basis
\[
 \{s_D:D\in\mathfrak N(P)\}.
\]
By \cref{cor:rigid-normal-basis}, the target has basis
\[
 \{\overline\eps_D:D\in\mathfrak N(P)\}.
\]
The map $\Xi$ sends the first basis element to the corresponding second basis
element.  It is therefore injective and hence an isomorphism.

Compose the quotient map
$\Hred\twoheadrightarrow\Hred/\Jsk$, the inverse of $\Xi$, and the
skein--cluster isomorphism $\Psi_{\rm sc}$ of
\cref{thm:skein-cluster}.  This gives \eqref{eq:global-Hall-cluster-map} and
identifies its kernel with $\Jsk$.  The statements on diagonals and
non-crossing direct sums follow from the definitions and
\eqref{eq:noncrossing-theta}.
\end{proof}

\subsection{Why this uses the entire Hall algebra}

Let $\Hall^{\rm red}_{\rm all}$ denote the central localization and
specialization of the entire cocycle-twisted exact Hall algebra, before any
restriction to rigid classes.  In the present notation this is exactly
$\Hred$.

\begin{corollary}[Whole-algebra statement]
\label{cor:whole-Hall}
The source of $\Phi_{\rm sk}$ contains every Hall basis object of
$\Mor(\proj A)$ after contractible central reduction.  In particular:
\begin{enumerate}[label=(\roman*),leftmargin=2.3em]
\item no composition subalgebra generated only by an initial cluster is used;
\item no non-rigid basis element is removed before taking the quotient;
\item the image of every product of arbitrary Hall elements is defined and is
      their product in the quantum cluster algebra.
\end{enumerate}
\end{corollary}

\begin{proof}
Every reduced indecomposable class is a diagonal generator, and
\cref{lem:indecomposable-generation} shows that those classes generate the
whole centrally reduced Hall algebra.  The map in
\cref{thm:Hall-skein-subquotient} is an algebra homomorphism on that algebra.
\end{proof}

Thus the construction is directly analogous in form to a Hall subquotient
realization:
\begin{equation}
\begin{CD}
 \Hall(\Mor(\proj A))
 @>{\text{bicharacter twist and localization}}>>
 \Hall_S^\star(\MM)[K_i^{-1}]\\
 @. @VV{\text{central coefficient specialization}}V\\
 @. \Hred @>{\text{quotient by }\Jsk}>> \Aq.
\end{CD}
\label{eq:subquotient-chain}
\end{equation}
A bicharacter twist is a standard Hall normalization rather than a quotient;
the genuine quotient is the explicit two-sided ideal $\Jsk$.

\subsection{Compatible and incompatible products}

For compatible rigid objects, the quotient does nothing new.

\begin{proposition}[Recovery of the chamber maps]
\label{prop:recover-chambers}
If $X$ and $Y$ are rigid and belong to one cluster, then
\begin{equation}
 \Phi_{\rm sk}(\eps_X\star_S\eps_Y)
 =\Phi_{\rm sk}(\eps_X)\Phi_{\rm sk}(\eps_Y)
 =v^{\Lambda_R(g_X,g_Y)}\Theta_{g(X\oplus Y)}.
 \label{eq:compatible-global-map}
\end{equation}
Hence $\Phi_{\rm sk}$ restricts to the Hall--theta chamber isomorphism of the
preceding paper.
\end{proposition}

\begin{proof}
The Hall product is the split direct sum with the quantum-torus exponent by
\eqref{eq:compatible-Hall-quantum}.  The image of the direct sum is the
normalized cluster monomial by \cref{thm:Hall-skein-subquotient}.
\end{proof}

The new content is the incompatible case.

\begin{theorem}[Global incompatible-product formula]
\label{thm:incompatible-product}
For arbitrary $x,y\in\Hred$,
\begin{equation}
 \Phi_{\rm sk}(x\star_S y)
 =\Phi_{\rm sk}(x)\Phi_{\rm sk}(y).
 \label{eq:global-multiplicativity}
\end{equation}
In particular, let $X,Y$ be incompatible rigid objects and expand their Hall
product as
\begin{equation}
 \eps_X\star_S\eps_Y
 =\sum_E h_{X,Y}^E(\omega)\eps_E.
 \label{eq:incompatible-Hall-expansion}
\end{equation}
Then
\begin{equation}
 \Theta_{g(X)}\Theta_{g(Y)}
 =\sum_E h_{X,Y}^E(\omega)
       \Phi_{\rm sk}(\eps_E),
 \label{eq:Hall-to-theta-structure}
\end{equation}
and every possibly non-rigid term $\Phi_{\rm sk}(\eps_E)$ has the finite
rigid theta expansion of \cref{sec:nonrigid}.
\end{theorem}

\begin{proof}
Equation \eqref{eq:global-multiplicativity} is the defining property of the
algebra homomorphism in \cref{thm:Hall-skein-subquotient}.  Apply it to the
basis expansion \eqref{eq:incompatible-Hall-expansion}.
\end{proof}

This theorem turns all exact Hall structure constants into quantum cluster
structure constants after one fixed quotient.  It does not assert that an
individual non-rigid Hall class is itself a theta basis element; generally it
is a linear combination of them.

\subsection{Coefficient specializations}

The most uniform theorem uses the frozen-localized Hall-compatible coefficient
system.  Several familiar forms follow by specialization.

\begin{corollary}[Coefficient regimes]
\label{cor:coefficient-regimes}
\begin{enumerate}[label=(\roman*),leftmargin=2.3em]
\item For arbitrary $n$, the canonical principal/walled quantization gives
      the isomorphism \eqref{eq:global-subquotient} after frozen localization.
\item If the mutable exchange matrix is full rank and
      $B_0^T\Lambda_0=dI$, one may take the global coefficient-free Hall
      cocycle associated with $\Lambda_0$; the quotient realizes the
      corresponding coefficient-free quantum cluster algebra after the chosen
      boundary specialization.
\item Setting $\omega=1$ gives a commutative quotient realization of the
      classical cluster algebra, with Ptolemy relations
      $x_\gamma x_\delta=x_{\gamma\#_0\delta}+x_{\gamma\#_\infty\delta}$.
\end{enumerate}
\end{corollary}

\begin{proof}
The first two statements use the globalizable Hall quantizations constructed
in the preceding paper and the coefficient skein models of
\cite{IshibashiKanoYuasa}.  The third is the classical specialization of the
skein and Hall relations.
\end{proof}

\subsection{Comparison with the acyclic morphism-Hall realization}

Ding--Xu--Zhang and Fu--Peng--Zhang obtain a principal-coefficient quantum
cluster algebra as a subquotient of the morphism Hall algebra of an acyclic
path algebra \cite{DingXuZhang,FuPengZhang}.  The present theorem shares three
features with that construction:
\begin{enumerate}[label=(\alph*),leftmargin=2.3em]
\item the source is an exact Hall algebra of projective morphisms;
\item contractible/projective directions are localized and specialized to
      coefficients;
\item the final cluster algebra is obtained by a quotient of an algebra
      generated by the whole relevant Hall category.
\end{enumerate}
The difference is the final ideal.  In the hereditary case the integration
map and Euler form already organize extension terms into the quantum
Caldero--Chapoton character.  In the nonhereditary type-$A$ case, the exact
integration is differential-blind; the explicitly generated Hall--skein ideal
supplies the missing opposite $2$-Calabi--Yau smoothing and straightens every
non-rigid term.
